%% ============================================================
%% Appendix — Software package, test suite, and performance
%% ============================================================

\appendix

% -------------------------------------------------------------------
\section{Software Package: \texttt{a6cw}}
\label{app:package}

All analysis code is organised into the \texttt{a6cw} Python package (version
1.0.0), hosted at \url{https://github.com/AlexBM173/a6-cw} with documentation
at \url{https://a6-cw.readthedocs.io}. The package separates reusable
algorithmic code from the two Jupyter notebooks, which contain only top-level
orchestration. It is installed in development mode via \texttt{pyproject.toml}
(\texttt{setuptools} build backend) and comprises two submodules:

\begin{center}
\begin{tabular}{ll}
\hline
\textbf{Module} & \textbf{Purpose} \\
\hline
\texttt{a6cw.ellipticity} & Galaxy shape estimators and response calibration \\
\texttt{a6cw.shear}       & Shear profiles, NFW theory, posterior inference \\
\hline
\end{tabular}
\end{center}

% -------------------------------------------------------------------
\subsection{\texttt{a6cw.ellipticity}}
\label{app:ellipticity}

This module implements three galaxy shape estimators and the
simulation-based calibration scheme for the Gaussian-weighted moment
estimator.

\subsubsection*{Ellipticity convention}

All three estimators return ellipticities in the \emph{distortion} convention.
The quadrupole moments of the surface brightness $I(x,y)$ are
\begin{equation}
    Q_{ij} = \frac{\sum_{pq} I(p,q)\,(p_i - \bar{p}_i)(p_j - \bar{p}_j)}
                  {\sum_{pq} I(p,q)},
\end{equation}
with flux-weighted centroid $\bar{p}_i$. The two ellipticity components are
\begin{equation}
    \epsilon_1 = \frac{Q_{11} - Q_{22}}{D}, \qquad
    \epsilon_2 = \frac{2\,Q_{12}}{D}, \qquad
    D = Q_{11} + Q_{22} + 2\sqrt{Q_{11}Q_{22} - Q_{12}^2}.
\end{equation}
For a pure ellipse with semi-axes $a \ge b$ this gives $|\epsilon| = (a-b)/(a+b)$,
with $\epsilon_1 > 0$ ($\epsilon_2 = 0$) for an east--west elongation.

\subsubsection*{Estimators}

\begin{description}
  \item[\texttt{unweighted\_ellipticity(image)}]
    Computes $(\epsilon_1, \epsilon_2)$ from the raw unweighted second moments
    of the pixel array. The centroid is found from unweighted first moments.
    Returns \texttt{(nan, nan)} when the total flux is non-positive, the
    determinant $Q_{11}Q_{22} - Q_{12}^2 < 0$ (noise-dominated), or the
    resulting ellipticity magnitude exceeds unity.

  \item[\texttt{weighted\_ellipticity\_raw(image, fwhm\_px, max\_iter, tol)}]
    Computes $(\epsilon_1, \epsilon_2)$ from intensity weighted by a circular
    Gaussian
    \begin{equation}
        W(x, y) = \exp\!\left(-\frac{(x-\bar{x}^W)^2 + (y-\bar{y}^W)^2}
                                     {2\sigma_w^2}\right), \quad
        \sigma_w = \frac{\mathrm{FWHM}}{2\sqrt{2\ln 2}}.
    \end{equation}
    The weighted centroid $(\bar{x}^W, \bar{y}^W)$ is found iteratively,
    re-centring the weight on the current estimate until the shift falls below
    \texttt{tol} pixels. The raw ellipticities returned are biased relative to
    the true reduced shear because the fixed circular weight suppresses the
    tangential signal; \texttt{apply\_simulation\_calibration} corrects for this.

  \item[\texttt{hsm\_ellipticity(image)}]
    Wraps \textsc{GalSim}'s \texttt{FindAdaptiveMom}
    \citep{Hirata2003, Mandelbaum2005}, which iteratively adapts an elliptical
    Gaussian weight to match the galaxy shape. Returns the reduced shear
    $(g_1, g_2)$ and the best-fit circular Gaussian size $\sigma_g$ in pixels.
\end{description}

\subsubsection*{Calibration}

The fixed circular weight introduces a multiplicative shear-response bias:
the measured ellipticity $\hat\epsilon$ is related to the true reduced shear
$g$ by $\hat\epsilon \approx R\,g$, where $R < 1$ depends on galaxy size
relative to the weight width. \texttt{calibrate\_weighted\_response} estimates
$R(\sigma_g)$ by drawing one noiseless \textsc{GalSim} \texttt{Exponential}
profile at each of twelve half-light radii, applying a known shear
$g_1 = 0.05$, and computing $R = \hat\epsilon_1 / g_1$. The result is a
look-up table; \texttt{apply\_simulation\_calibration} divides each galaxy's
raw ellipticities by its interpolated $R(\sigma_{g,i})$ using
\texttt{np.interp}, clamping at the boundary of the calibrated range.

\subsubsection*{Batch processing}

\texttt{measure\_all\_ellipticities(stamps\_path)} loads all galaxy stamps
from a multi-extension FITS file and applies all three estimators in parallel
using a \texttt{ThreadPoolExecutor}. The Numba JIT kernels (see
Section~\ref{app:performance}) and \textsc{GalSim}'s C\texttt{++} HSM
implementation both release the Python GIL, enabling genuine thread-level
parallelism.

% -------------------------------------------------------------------
\subsection{\texttt{a6cw.shear}}
\label{app:shear}

This module implements the full weak-lensing shear pipeline from galaxy
positions and ellipticities to a Bayesian posterior on halo mass and source
redshift.

\subsubsection*{Module-level constants}

Radial binning is fixed at module import time to ensure consistency across all
pipeline functions:
\begin{align*}
    &\texttt{N\_BINS} = 10, \\
    &\texttt{BIN\_EDGES} = \mathrm{logspace}(\log_{10}100,\, \log_{10}3600,\, 11)
      \text{ [arcsec]}, \\
    &\texttt{BIN\_CENTRES}_k = \sqrt{\texttt{BIN\_EDGES}_k \cdot \texttt{BIN\_EDGES}_{k+1}}
      \quad\text{(geometric mean)}.
\end{align*}

\subsubsection*{Pipeline functions}

\begin{description}
  \item[\texttt{load\_positions(path)}]
    Reads a two-column whitespace-separated text file of $(x, y)$ positions
    in arcseconds, returning an \texttt{(N, 2)} array.

  \item[\texttt{compute\_tangential\_ellipticities(halo\_pos, source\_pos,
        $\epsilon_1$, $\epsilon_2$)}]
    Forms all $N_h \times N_s$ halo--source pairs by NumPy broadcasting.
    The tangential and cross components are computed via trig-free double-angle
    identities,
    \begin{equation}
        \cos 2\phi = \frac{dx^2 - dy^2}{r^2}, \qquad
        \sin 2\phi = \frac{2\,dx\,dy}{r^2},
    \end{equation}
    eliminating \texttt{arctan2}, \texttt{cos}, and \texttt{sin} calls over the
    full pair array. Returns three flat arrays of length $N_h \times N_s$
    containing the separation $r$ [arcsec], the tangential component
    $\epsilon_t$ (E-mode), and the cross component $\epsilon_\times$ (B-mode).

  \item[\texttt{mean\_tangential\_shear(r\_pairs, $\epsilon_t$, bin\_edges, weights)}]
    Bins tangential ellipticities into the ten logarithmically-spaced radial
    bins. In the weak-lensing regime $\langle\epsilon_t\rangle = \gamma_t$, so
    the mean tangential ellipticity directly estimates the shear.
    Bin assignment uses a single \texttt{np.digitize} pass [$O(N\log B)$]; per-bin
    sums use \texttt{np.bincount}. The unweighted uncertainty is the
    Bessel-corrected standard error on the mean; an optional
    inverse-variance weighted path is also provided. Bins with fewer than two
    valid pairs return \texttt{nan}.

  \item[\texttt{build\_nfw\_theory(mass, source\_redshift, ...)}]
    Evaluates the predicted mean tangential shear at \texttt{BIN\_CENTRES} for
    an NFW halo using the \texttt{NFWHalo\_theory} helper (a thin wrapper
    around \textsc{GalSim}'s \texttt{NFWHalo}). Default parameters
    ($c = 4$, $z_l = 0.3$, $\Omega_m = 0.3$, $\Omega_\Lambda = 0.7$) match
    the simulation truth.

  \item[\texttt{log\_likelihood(gt\_data, gt\_err, mass, source\_redshift)}]
    Evaluates the Gaussian log-likelihood
    \begin{equation}
        \ln\mathcal{L} = -\frac{1}{2}
        \sum_k \left(\frac{\hat\gamma_t^k - \gamma_t^{\mathrm{th}}(k)}
                          {\sigma_k}\right)^2,
    \end{equation}
    summing only over bins where both data and theory are finite. Returns
    $-\infty$ if no valid bins remain.

  \item[\texttt{compute\_posterior\_grid(gt\_data, gt\_err, mass\_grid,
        zs\_grid)}]
    Evaluates the unnormalised log-posterior on a 2-D $(M,\,z_s)$ grid under
    flat priors, distributing cells across OS processes via
    \texttt{ProcessPoolExecutor} (with a serial fallback for grids smaller than
    50 cells). Returns an $(N_M, N_{z_s})$ array.
\end{description}

\subsubsection*{Plotting utilities}

The module also provides \texttt{plot\_tangential\_shear\_profiles},
\texttt{plot\_cross\_shear\_null\_test}, \texttt{plot\_joint\_posterior},
\texttt{plot\_mass\_posterior\_fixed\_zs}, and \texttt{plot\_best\_fit\_overlay},
all writing PNG files to the \texttt{outputs/} directory.
\texttt{plot\_joint\_posterior} and \texttt{plot\_mass\_posterior\_fixed\_zs}
compute credible-interval boundaries using
\texttt{scipy.integrate.cumulative\_trapezoid}, which performs the CDF
integration in a single $O(N)$ pass.

% -------------------------------------------------------------------
\section{Test Suite}
\label{app:tests}

The package is tested by 82 unit tests spread across two files under
\texttt{tests/}, runnable with
\begin{verbatim}
python -m pytest tests/ -v
\end{verbatim}
Tests run automatically on every push via a GitHub Actions workflow
(\texttt{.github/workflows/tests.yml}) using Python~3.11 on
\texttt{ubuntu-latest}. No files from \texttt{data/} or \texttt{outputs/} are
required; all tests complete in under five seconds.

\subsection*{Design principles}

\textbf{Analytic ground truth where possible.} Most tests construct inputs with
known mathematical properties and verify the output against a closed-form
prediction, rather than comparing to a stored baseline. A failing test
therefore points directly to the formula that is wrong.

\textbf{\textsc{GalSim} profiles for shape estimators.} The three estimators
operate on \texttt{galsim.Image} objects. Tests draw simple noiseless profiles
(Gaussians with and without applied shear) using \texttt{drawImage} with
\texttt{method='no\_pixel'}, which gives a pixel array matching the analytical
profile to machine precision.

\textbf{Edge cases alongside the happy path.} Every estimator test includes a
zero-flux (or empty) input that should return \texttt{nan}, alongside the
positive-signal cases.

\textbf{Separation of concerns.} \texttt{test\_ellipticity.py} covers only
\texttt{a6cw.ellipticity}; \texttt{test\_shear.py} covers only
\texttt{a6cw.shear}, making failures easy to localise.

\subsection{\texttt{tests/test\_ellipticity.py} (39 tests)}

\begin{description}
  \item[\texttt{TestPixelGrids} (6 tests)]
    Verifies that \texttt{\_pixel\_grids} returns arrays of the correct shape,
    that grids sum to zero (symmetric about the centre), that the centre pixel
    is exactly zero for an odd stamp, and that $x$ varies along columns while
    $y$ varies along rows.

  \item[\texttt{TestGaussianWeight} (4 tests)]
    Confirms the circular Gaussian weight $W = e^{-(dx^2+dy^2)/2\sigma^2}$ peaks
    at unity at the origin, falls to $e^{-1/2}$ at one sigma, is symmetric in
    $x$ and $y$, and is strictly positive everywhere.

  \item[\texttt{TestUnweightedEllipticity} (8 tests)]
    Checks that a circular Gaussian galaxy gives $(\epsilon_1,\epsilon_2)
    \approx (0,0)$ to $<10^{-6}$; that a zero-flux stamp returns
    \texttt{(nan,\,nan)}; that axis-aligned sheared Gaussians recover the
    closed-form value $\epsilon_1 = (\sigma_x - \sigma_y)/(\sigma_x + \sigma_y)$
    to $0.1\%$; that a 45$^\circ$-rotated galaxy produces pure $\epsilon_2$; that
    a \textsc{GalSim} Gaussian with $g_1 = 0.3$ returns $\epsilon_1 = 0.3$ to
    $<10^{-5}$ relative error (the exact identity $\epsilon_1 = g_1$ for a
    Gaussian); and that results are bounded by unity.

  \item[\texttt{TestWeightedEllipticity} (6 tests)]
    Analogous tests for the Gaussian-weighted estimator, including a check that
    a wider weight FWHM increases the raw $\epsilon_1$ response (because it
    captures more of the galaxy's ellipticity signal).

  \item[\texttt{TestHSMEllipticity} (6 tests)]
    Verifies that circular galaxies give zero HSM shear, that $\sigma_g > 0$,
    that $g_1$ and $g_2$ are recovered to $<10^{-4}$ relative error on noiseless
    Gaussian profiles, that empty images return \texttt{nan}, and that $\sigma_g$
    increases monotonically with galaxy size.

  \item[\texttt{TestOLSThroughOrigin} (4 tests)]
    Confirms the force-through-origin OLS regression recovers a known slope
    exactly, recovers a noisy slope to $<1\%$ with 2000 points, filters
    \texttt{nan} values, and recovers negative slopes.

  \item[\texttt{TestApplySimulationCalibration} (5 tests)]
    Verifies that unit response is the identity; half response doubles
    ellipticity; \texttt{nan} ellipticities remain \texttt{nan}; galaxies with
    \texttt{nan} $\sigma_g$ use the sample median before interpolation; and
    $\sigma_g$ values outside the calibration range are clamped to the nearest
    boundary node.
\end{description}

\subsection{\texttt{tests/test\_shear.py} (43 tests)}

\begin{description}
  \item[\texttt{TestModuleConstants} (7 tests)]
    Checks that \texttt{BIN\_EDGES} has $N_\mathrm{bins}+1$ elements,
    \texttt{BIN\_CENTRES} are the geometric means of adjacent edges, edges are
    strictly increasing, the range spans 100--3600 arcsec, and that
    \texttt{COLOURS} and \texttt{LABELS} both contain the three method keys.

  \item[\texttt{TestLoadPositions} (3 tests)]
    Verifies round-trip I/O: write with \texttt{np.savetxt}, read with
    \texttt{load\_positions}, confirm shape $(N, 2)$ and value preservation
    (including negative and decimal coordinates).

  \item[\texttt{TestComputeTangentialEllipticities} (7 tests)]
    Uses three exact analytic cases for the rotation formula:
    \begin{center}
    \begin{tabular}{llll}
    \hline
    Source direction & $\phi$ & $\epsilon_t$ & $\epsilon_\times$ \\
    \hline
    East ($x > 0$) & $0$ & $-\epsilon_1$ & $0$ \\
    North ($y > 0$) & $\pi/2$ & $+\epsilon_1$ & $0$ \\
    Northeast ($\phi = \pi/4$, $\epsilon_2 = 1$) & $\pi/4$ & $-\epsilon_2$ & $0$ \\
    \hline
    \end{tabular}
    \end{center}
    Also checks the 3--4--5 right-triangle separation formula, output shape for
    $N_h = 7$, $N_s = 13$, and that zero input ellipticity gives zero output.

  \item[\texttt{TestMeanTangentialShear} (7 tests)]
    Uses a two-bin layout with four known pairs to test: correct mean, pair
    count, standard error ($\sigma = \sqrt{0.02}/\sqrt{2} = 0.1$ for
    $\{0.4,\,0.6\}$); \texttt{nan} for single-pair bins; \texttt{nan} exclusion
    from counts; empty-bin \texttt{nan}; and zero error for uniform values.

  \item[\texttt{TestEstimateShapeNoiseVariance} (4 tests)]
    Confirms known variance ($\{1,{-}1\}$ gives 2 with \texttt{ddof=1}),
    \texttt{nan} exclusion, zero variance for constant arrays, and a
    \texttt{float} return type.

  \item[\texttt{TestOptimalWeights} (4 tests)]
    Checks equal weights for uniform noise, zero weight for \texttt{nan} pairs,
    positive weight for finite pairs, and $w = 1/\hat\sigma_\epsilon^2$.

  \item[\texttt{TestBuildNFWTheory} (6 tests)]
    Verifies output length equals \texttt{N\_BINS}; all values positive; profile
    decreasing with radius; higher mass gives stronger shear at all radii; higher
    source redshift gives stronger shear; and non-default cosmological parameters
    are accepted.

  \item[\texttt{TestLogLikelihood} (5 tests)]
    Confirms a perfect fit gives $\ln\mathcal{L} = 0$; an imperfect fit gives
    $\ln\mathcal{L} < 0$; a larger residual gives a more negative
    log-likelihood; \texttt{nan} bins are silently excluded; and all-\texttt{nan}
    data returns $-\infty$.
\end{description}

% -------------------------------------------------------------------
\section{Performance and Profiling}
\label{app:performance}

All benchmarks were measured on WSL2 (Windows Subsystem for Linux),
Python~3.11.11, NumPy~2.x, \textsc{GalSim}~2.8.4, and Numba~0.65.1.
Wall-clock times are medians over 50--100 repeated calls with one warm-up call
excluded. Peak heap memory is measured with Python's \texttt{tracemalloc}.
Profiling code and plot generation are in \texttt{profiling/run\_profiling.py};
plots are written to \texttt{outputs/profiling/}.

\subsection{Optimisations implemented}

Eleven optimisations were applied across the two modules and the vendored
\texttt{nfw\_theory.py} file.

\begin{enumerate}
  \item \textbf{Numba JIT for unweighted moments.}
    The three-pass pixel loop in \texttt{unweighted\_ellipticity} is compiled by
    Numba (\texttt{@njit(cache=True, nogil=True)}). Intermediate quantities are
    accumulated in scalar registers, eliminating all temporary NumPy array
    allocations.

  \item \textbf{Numba JIT for weighted moments.}
    The same treatment for \texttt{weighted\_ellipticity\_raw}, with additional
    micro-optimisations: \texttt{inv\_2s2} $= 1/(2\sigma^2)$ is precomputed
    once outside the centroid loop, and convergence is tested as
    $|\Delta\bar{x}|^2 < \epsilon^2$ to avoid a square root per iteration.

  \item \textbf{\texttt{lru\_cache} on pixel grids.}
    The coordinate offset grids are cached by stamp shape via
    \texttt{functools.lru\_cache(maxsize=32)}, so repeated calls on same-size
    stamps incur no allocation cost.

  \item \textbf{\texttt{ThreadPoolExecutor} in \texttt{measure\_all\_ellipticities}.}
    Both Numba JIT kernels and \textsc{GalSim}'s C\texttt{++} HSM implementation
    release the Python GIL, so a thread pool achieves genuine parallelism across
    galaxy stamps.

  \item \textbf{Single \textsc{GalSim} draw per HLR in
        \texttt{calibrate\_weighted\_response}.}
    Noiseless stamps with fixed parameters are fully deterministic; the original
    loop drew 1000 pixel-identical stamps and averaged them. Reducing to one draw
    per HLR gives a ${\sim}1000\times$ reduction in GalSim renders
    ($12\,000 \to 12$).

  \item \textbf{Trig-free double-angle identities in
        \texttt{compute\_tangential\_ellipticities}.}
    $\cos 2\phi = (dx^2-dy^2)/r^2$ and $\sin 2\phi = 2\,dx\,dy/r^2$ eliminate
    \texttt{arctan2}, \texttt{cos}, and \texttt{sin} calls over the
    $N_h \times N_s$ pair array.

  \item \textbf{\texttt{np.digitize\,+\,np.bincount} in
        \texttt{mean\_tangential\_shear} (unweighted).}
    A single $O(N\log B)$ digitize pass replaces $B$ repeated boolean
    masks [$O(NB)$ total].

  \item \textbf{\texttt{np.bincount} weighted path.}
    Per-bin loops for the inverse-variance weighted sums are replaced with
    vectorised \texttt{np.bincount} calls.

  \item \textbf{\texttt{ProcessPoolExecutor} in \texttt{compute\_posterior\_grid}.}
    The NFW theory evaluation holds the GIL; separate OS processes bypass it.

  \item \textbf{\texttt{cumulative\_trapezoid} for CDF computation.}
    The $O(N^2)$ list comprehension
    \texttt{[np.trapezoid(p[:k+1], x[:k+1]) for k in range(N)]} is replaced
    by \texttt{scipy.integrate.cumulative\_trapezoid} [$O(N)$], which also
    avoids accumulation of floating-point rounding from repeated prefix sums.

  \item \textbf{Vectorised \texttt{NFWHalo.\_\_gamma} in
        \texttt{nfw\_theory.get\_tangential\_shear}.}
    \textsc{GalSim}'s private \texttt{\_\_gamma(r, ks)} accepts arrays;
    broadcasting \texttt{ks\_arr = np.full\_like(rs, ks)} replaces a Python
    loop over 10 radial bins with one vectorised C-level call.
\end{enumerate}

\subsection{Results}

Table~\ref{tab:perf} summarises the before/after wall-clock times at
representative input sizes, and Table~\ref{tab:mem} summarises peak heap
memory for the two moment estimators.

\begin{table}[h]
\centering
\caption{Wall-clock performance before and after all optimisations (medians).
         The \texttt{calibrate\_weighted\_response} entry combines optimisations
         5 and 4; the posterior grid entry combines optimisations 9 and 11.}
\label{tab:perf}
\small
\begin{tabular}{lrrr}
\hline
\textbf{Routine (input size)} & \textbf{Before} & \textbf{After} & \textbf{Speedup} \\
\hline
\texttt{unweighted\_ellipticity} ($48\times48$ px) & 0.030\,ms & 0.007\,ms & $4\times$ \\
\texttt{weighted\_ellipticity\_raw} ($48\times48$ px) & 0.066\,ms & 0.025\,ms & $2.6\times$ \\
\texttt{hsm\_ellipticity} ($48\times48$ px) & 0.184\,ms & 0.175\,ms & $1.1\times$ \\
\texttt{calibrate\_weighted\_response} (12 HLRs) & ${\sim}14\,\mathrm{s}$ & 60\,ms & ${\sim}230\times$ \\
\texttt{compute\_tangential\_ellipticities} ($10^5$ sources) & 84\,ms & 39\,ms & $2.2\times$ \\
\texttt{build\_nfw\_theory} (10 bins) & 0.421\,ms & 0.070\,ms & $6\times$ \\
\texttt{compute\_posterior\_grid} ($2\,500$ cells) & 1\,217\,ms & 131\,ms & $9.3\times$ \\
\hline
\end{tabular}
\end{table}

\begin{table}[h]
\centering
\caption{Peak heap memory for the two pixel-based moment estimators at
         $48\times48$ pixels. HSM (C\texttt{++}) is unchanged.}
\label{tab:mem}
\begin{tabular}{lrrr}
\hline
\textbf{Estimator} & \textbf{Before} & \textbf{After} & \textbf{Reduction} \\
\hline
Unweighted moments & 127\,KB & 18\,KB & $7\times$ \\
Gaussian-weighted moments & 181\,KB & 18\,KB & $10\times$ \\
HSM (\textsc{GalSim}) & 13\,KB & 13\,KB & unchanged \\
\hline
\end{tabular}
\end{table}

The 18\,KB residual after optimisation equals the input array itself
($48 \times 48 \times 8\,\mathrm{bytes} = 18\,\mathrm{KB}$), confirming that
the Numba JIT kernels allocate no intermediate arrays.

The dominant gain is \texttt{calibrate\_weighted\_response}: the ${\sim}230\times$
speedup (not $1000\times$) reflects fixed per-HLR overhead in \textsc{GalSim}
cosmology initialisation and pixel convolution that is unaffected by the draw-count
reduction. The posterior-grid speedup of $9.3\times$ decomposes into approximately
$6\times$ from the vectorised NFW call (opt.~11) and $1.6\times$ additional from
\texttt{ProcessPoolExecutor} (opt.~9) at 2\,500 cells. The coursework grid
($100 \times 50 = 5\,000$ cells) completes in ${\approx}0.35\,\mathrm{s}$
versus the original ${\approx}2.5\,\mathrm{s}$ ($7\times$).

\subsection{Scaling behaviour}

The ellipticity estimators scale as $O(N_\mathrm{px})$ in the number of pixels
($N_\mathrm{px} = N^2$ for an $N\times N$ stamp). At small stamps the measured
log--log slope is below unity (0.84 for unweighted, 1.12 for weighted) because
fixed Numba dispatch overhead dominates at a few hundred pixels; above
$128\times128$ all estimators approach the asymptotic linear regime.

\texttt{compute\_tangential\_ellipticities} scales as $O(N_h \cdot N_s)$, with
a measured log--log slope of ${\approx}1.0$ confirming linear scaling with
source count at fixed halo count. At $10^5$ sources and 10 halos, peak
memory is ${\approx}77\,\mathrm{MB}$ ($10^6$ double-precision pairs $\times$
three arrays).

The posterior grid scales as $O(N_M \cdot N_{z_s} / N_\mathrm{cores})$ with
\texttt{ProcessPoolExecutor}. At 2\,500 cells the per-cell amortised cost is
0.05\,ms, compared to 0.49\,ms in the original serial implementation.
